% from ega4-18-fr.tex
\oldpage[IV-4]{200}

voisinage ouvert affine dont l'anneau est de la forme $C_{(t)}$, où $t$ est un élément homogène de degré $1$ de $C$; comme par hypothèse $C$ est un $A$-module plat, il en est de même de son anneau de fractions $C_t$ et de la composante de degré $0$, $C_{(t)}$, de ce $A$-module gradué, d'où notre assertion.
  \item[(ii)] On peut toujours se borner au cas où $S$ et $X$ sont affines, $B$ étant donc une $A$-algèbre de présentation finie, $\mathfrak J$ un idéal de type fini de $B$. En vertu de (8.9.1), (8.6.3) et (11.2.9), il existe un sous-anneau noethérien $A_0$ de $A$, une $A_0$-algèbre de type fini $B_0$, un idéal $\mathfrak J_0$ de $B_0$ tels que $B=B_0\otimes_{A_0}A$, $\mathfrak J=\mathfrak J_0B$, que $\operatorname{gr}_{\mathfrak J_0}^*(B_0)$ soit un $A_0$-module \emph{plat} et que $\operatorname{gr}_{\mathfrak J}^*(B)=\operatorname{gr}_{\mathfrak J_0}^*(B_0)\otimes_{A_0}A$. En vertu de (i), on a alors $\bigoplus_{n\geq0}\mathfrak J^n=(\bigoplus_{n\geq0}\mathfrak J_0^n)\otimes_{A_0}A=(\bigoplus_{n\geq0}\mathfrak J_0^n)\otimes_{B_0}B$; comme $\bigoplus_{n\geq0}\mathfrak J_0^n$ est une $B_0$-algèbre engendrée par l'idéal de type fini $\mathfrak J_0$ de $B_0$, donc une $B_0$-algèbre de type fini, et par suite de présentation finie puisque $B_0$ est noethérien, on en déduit que $\bigoplus_{n\geq0}\mathfrak J^n$ est une $B$-algèbre de présentation finie. De même, le morphisme $\operatorname{Proj}(\bigoplus_{n\geq0}\mathfrak J_0^n)\to\Spec(B_0)$ est de type fini (II, 2.7.1), et par suite de présentation finie; donc le morphisme $X'\to X$, qui s'en déduit par changement de base, est de présentation finie. On sait déjà (19.4.2) que l'immersion canonique $Y'\to X'$ est régulière; comme $Y'$ est $S$-plat, ainsi qu'on l'a vu dans (i), on déduit de (19.2.4) que l'immersion $Y'\to X'$ est transversalement régulière relativement à $S$ et que $X'$ est $S$-plat aux points de $Y'$, donc (11.3.1) dans un voisinage ouvert de $Y'$. D'autre part, le fait que $\operatorname{gr}_{\mathscr J}^*(\sh{O}_X)$ est une $\sh{O}_X$-Algèbre $S$-plate entraîne, par (11.3.4), que $X$ est $S$-plat aux points de $Y$, donc (11.3.1) dans un voisinage de $Y$.
\end{enumerate}

\begin{corollary}[19.4.7]
\label{IV.19.4.7}
Soient $g:X\to S$ un morphisme localement de présentation finie, $Y$ un sous-préschéma fermé de $X$ tel que le morphisme composé $Y\to X\to S$ soit localement de présentation finie; supposons en outre que $Y$ soit $S$-plat et que $\sh{O}_X$ soit normalement plat le long de $Y$ (11.3.4). Alors (avec les notations de (19.4.1)), le morphisme $X'\to X$ est de présentation finie, l'immersion canonique $Y'\to X'$ est transversalement régulière de codimension $1$ relativement à $S$, $Y'$ est $Y$-plat (donc $S$-plat) et l'ensemble des points de $X$ (resp. $X'$) où $X$ (resp. $X'$) est $S$-plat est un voisinage de $Y$ (resp. $Y'$).
\end{corollary}

En effet, par hypothèse $\operatorname{gr}_{\mathscr J}^*(\sh{O}_X)$ est un $(\sh{O}_X/\mathscr J)$-Module plat, donc est $S$-plat puisque $Y$ est supposé être $S$-plat. On peut par suite appliquer les résultats de (19.4.6, (i) et (ii)). Comme $Y'$ est $Y$-isomorphe à $\operatorname{Proj}(\operatorname{gr}_{\mathscr J}^*(\sh{O}_X))$ (19.4.2), le même raisonnement que dans (19.4.6) montre que $Y'$ est $Y$-plat.

\begin{proposition}[19.4.8]
\label{IV.19.4.8}
Supposons vérifiées les hypothèses de (19.4.7) et en outre que le morphisme $Y'\to Y$ soit lisse. Soient $x'$ un point de $Y'$, $x$ son image dans $Y$. Pour que $Y$ soit lisse sur $S$ au point $x$, il faut et il suffit que $Y'$ soit lisse sur $S$ au point $x'$; lorsqu'il en est ainsi, $X'$ est alors lisse sur $S$ au point $x'$.

Les conclusions précédentes sont en particulier vérifiées lorsque le morphisme $X\to S$ est localement de présentation finie et que l'immersion $Y\to X$ est transversalement régulière relativement à $S$.
\end{proposition}

Compte tenu de la compatibilité de la formation d'un préschéma éclaté avec les changements de base dans le cas envisagé (19.4.6, (i)) et de (17.5.1) et (17.7.1),

